\thispagestyle{abstractpagestyle}

\vspace*{36pt}

\begin{center}
    \textbf{\fontsize{20pt}{24pt} \selectfont ABSTRACT}
\end{center}

\vspace{24pt}

Quantum image processing promises accelerations for filtering and edge detection, yet noisy intermediate-scale quantum (NISQ) devices impose shallow circuits, limited qubit counts, and significant noise. The flexible representation of quantum images (FRQI) encodes pixel intensities as rotation angles steered by multi-controlled addresses; naive preparation using multi-controlled $X$ gates without dedicated ancillas yields large transpiled depth and controlled-NOT (CX) counts.

This dissertation thesis compares structural FRQI state preparation with a \emph{no-ancilla} multi-controlled decomposition against a reusable \emph{v-chain} ancilla pattern on $4\times 4$, $8\times 8$, and $16\times 16$ test images using the Qiskit Aer simulator. Reports consist of transpiled resources (qubits, depth, CX), noisy FRQI reconstructions under synthetic depolarizing noise, quantum Hadamard edge detector (QHED) edge metrics against classical Sobel references, and a proof-of-concept readout mitigation slice based on inverting an empirical $2\times 2$ confusion matrix on the color qubit.

The thesis body is organized into dedicated methods and results chapters so that resource accounting (optimization level~3 structural transpiles) can be distinguished from the frozen noisy bundles (optimization level~0) documented in the repository manifests.

At $16\times 16$, v-chain preparation reaches 11\,663 transpiled CX gates versus a linearly scaled naive estimate of 114\,176 CX from one representative slice; ideal FRQI reconstruction matches the classical images (SSIM $=1$), while QHED--Sobel similarity degrades at higher resolution (SSIM $\approx 0.16$ at $8\times 8$, $\approx 0.01$ at $16\times 16$). Paired naive--v-chain comparisons under noise for \texttt{edge\_ssim} are exploratory (Wilcoxon $p\approx 0.078$); readout mitigation on $4\times 4$ shows a mean shift in inferred $p_1$ with $p\approx 0.002$. All evidence is simulator-only (experiments haven't been executed on IBM backend); some $16\times 16$ v-chain experiments are absent when density-matrix simulation exceeds practical memory limits (insufficient RAM).

\vfill
